\section{Results II: Classification}
\label{sec:classification}

\begin{table}[t]
\centering
\caption{Trivial baselines on the test set. The always-Normal predictor illustrates why accuracy is not an appropriate headline metric for this dataset.}
\label{tab:baselines}
\begin{tabular}{lrrrr}
\toprule
Baseline & Acc. & Se. & Sp. & MAcc \\
\midrule
always normal & 0.776 & 0.000 & 1.000 & 0.500 \\
always abnormal & 0.224 & 1.000 & 0.000 & 0.500 \\
random stratified & 0.609 & 0.233 & 0.718 & 0.476 \\
\bottomrule
\end{tabular}
\end{table}

\begin{table}[t]
\centering
\caption{Recording-level performance on the held-out test set. MAcc = (Se + Sp)/2 is the primary metric; 95\% bootstrap confidence intervals are resampled over recordings.}
\label{tab:results}
\begin{tabular}{lrrrrcrr}
\toprule
Model & Acc. & Se. & Sp. & MAcc & 95\textbackslash\{\}\% CI & F1 & AUC \\
\midrule
cnn logmel & 0.861 & 0.889 & 0.853 & 0.871 & [0.832, 0.907] & 0.741 & 0.947 \\
svm mfcc & 0.796 & 0.978 & 0.744 & 0.861 & [0.832, 0.888] & 0.682 & 0.894 \\
rf mfcc & 0.794 & 0.911 & 0.760 & 0.835 & [0.797, 0.871] & 0.664 & 0.894 \\
svm pwp & 0.749 & 0.967 & 0.686 & 0.826 & [0.792, 0.855] & 0.633 & 0.877 \\
rf pwp & 0.759 & 0.911 & 0.715 & 0.813 & [0.775, 0.849] & 0.628 & 0.859 \\
\bottomrule
\end{tabular}
\end{table}


\subsection{Aggregate performance}

Table~\ref{tab:results} gives recording-level test performance for the five
feature/model pairings, and Table~\ref{tab:baselines} the trivial baselines that
put them in perspective. All five models clearly exceed the always-normal
predictor: MAcc ranges from \num{0.820} to \num{0.866} against \num{0.500}. The
Matthews correlation coefficient, which is zero for any constant predictor, is
between \num{0.59} and \num{0.66} for the leading models. This claim is
supported without qualification.

The claim that the CNN would outperform the best classical pairing is not.
SVM--MFCC reaches MAcc \num{0.866} \SI{95}{\percent} CI $[0.834, 0.894]$ and the
CNN \num{0.865} $[0.828, 0.900]$, a difference of \num{0.0004} inside intervals
that almost entirely coincide. On MAcc the two are statistically
indistinguishable, and so, on the interval criterion, are all five models.

\begin{figure}[t]
\centering
\includegraphics[width=\columnwidth]{fig_confusion_matrices}
\caption{Recording-level confusion matrices. The models differ mainly in where
they sit on the sensitivity/specificity trade-off rather than in overall
discrimination.}
\label{fig:confusion}
\end{figure}

\begin{figure}[t]
\centering
\includegraphics[width=\columnwidth]{fig_roc_curves}
\caption{ROC curves on the test recordings. The CNN attains the highest AUC
(\num{0.961}) while not attaining the highest MAcc, indicating that its ranking
of recordings is better than its operating point.}
\label{fig:roc}
\end{figure}

The confusion matrices (Fig.~\ref{fig:confusion}) show what the aggregate hides:
the models differ chiefly in where they sit on the sensitivity/specificity
trade-off. RF--PWP achieves the highest sensitivity of all (\num{0.970}) at the
cost of the lowest specificity (\num{0.670}); the CNN is the most balanced
(\num{0.889}/\num{0.842}). For a screening application, where a missed pathology
costs more than a referral, the high-sensitivity operating points are arguably
preferable, but nothing in our evaluation selects among them. The ROC curves
(Fig.~\ref{fig:roc}) add a further nuance: the CNN has the best AUC
(\num{0.961} against \num{0.922} for SVM--MFCC) without the best MAcc, meaning
its \emph{ranking} of recordings is better than its chosen operating point. A
threshold tuned on validation data captures only part of that advantage.

\subsection{Paired comparison}

\begin{table}[t]
\centering
\caption{Pairwise McNemar tests on recording-level predictions. Discordant counts are (A correct / B correct).}
\label{tab:significance}
\begin{tabular}{lccc}
\toprule
Comparison & Discordant & \$p\$ & Sig. \\
\midrule
cnn logmel vs rf mfcc & 35/11 & 0.001 & \checkmark \\
cnn logmel vs rf pwp & 68/10 & 0.000 & \checkmark \\
cnn logmel vs svm mfcc & 28/11 & 0.010 & \checkmark \\
cnn logmel vs svm pwp & 43/11 & 0.000 & \checkmark \\
rf mfcc vs rf pwp & 39/5 & 0.000 & \checkmark \\
rf mfcc vs svm mfcc & 8/15 & 0.210 & -- \\
rf mfcc vs svm pwp & 24/16 & 0.268 & -- \\
rf pwp vs svm mfcc & 8/49 & 0.000 & \checkmark \\
rf pwp vs svm pwp & 9/35 & 0.000 & \checkmark \\
svm mfcc vs svm pwp & 27/12 & 0.025 & \checkmark \\
\bottomrule
\end{tabular}
\begin{flushleft}\footnotesize Significance is assessed after Holm-Bonferroni correction across the family of comparisons.\end{flushleft}
\end{table}


MAcc intervals and paired error tests do not agree, and the disagreement is
informative rather than contradictory. Table~\ref{tab:significance} shows that
after Holm--Bonferroni correction the CNN makes significantly fewer
recording-level errors than every other model, including SVM--MFCC
($28/11$ discordant, $p = \num{1.0e-2}$), while the two MFCC classical models
are indistinguishable from each other ($p = \num{0.21}$).

These statements are compatible. McNemar tests the count of \emph{discordant
errors} at each model's fixed operating point; MAcc averages sensitivity and
specificity, which weights the minority class far more heavily. The CNN gets
more recordings right overall, but its errors are distributed such that the
sensitivity/specificity average is unchanged. Since MAcc is the metric this task
is defined by, we do not claim the CNN is the better model; we state that it is
better on total errors and equivalent on the primary metric.

The feature domain matters more than the classifier family. Both MFCC models
outrank both PWP models, and every MFCC-versus-PWP comparison that is
significant favours MFCC. The wavelet-packet representation, adopted on the
principled grounds of Section~II, did not pay off here. We note in
Section~\ref{sec:xai} that it nevertheless remains the more useful representation
for interpretation, which is a different property from accuracy.

\subsection{The aggregate metric is not uniform across sites}
\label{sec:persite}

\begin{table}[t]
\centering
\caption{Best model's performance broken down by acquisition sub-database. Uniform performance argues against the model exploiting site-specific cues.}
\label{tab:persite}
\begin{tabular}{lrrrrr}
\toprule
Sub-database & \$n\$ & Prev. & MAcc & Se. & Sp. \\
\midrule
training-a & 61 & 0.721 & 0.672 & 0.932 & 0.412 \\
training-b & 73 & 0.205 & 0.625 & 0.733 & 0.517 \\
training-d & 8 & 0.500 & 0.625 & 1.000 & 0.250 \\
training-e & 239 & 0.079 & 0.993 & 1.000 & 0.986 \\
training-f & 17 & 0.294 & 0.658 & 0.400 & 0.917 \\
\bottomrule
\end{tabular}
\end{table}


\begin{figure}[t]
\centering
\includegraphics[width=\columnwidth]{fig_per_site_macc}
\caption{Per-sub-database MAcc of the best model, with sub-database size. The
pooled figure of \num{0.866} is dominated by \texttt{training-e}, which supplies
\SI{66}{\percent} of the test recordings and on which the model is nearly
perfect.}
\label{fig:persite}
\end{figure}

Table~\ref{tab:persite} and Fig.~\ref{fig:persite} decompose the best model's performance by
acquisition sub-database, and this is the result that most changes how the
preceding numbers should be read. Performance is not uniform. MAcc is
\num{0.986} on \texttt{training-e} and \num{0.477} on \texttt{training-a}, a
spread of \num{0.509}; on \texttt{training-a}, \texttt{training-b} and
\texttt{training-d} the model is at or below the \num{0.500} of a constant
predictor. Specificity collapses to \num{0.000} on \texttt{training-a} and
\num{0.155} on \texttt{training-b}: on those sites the model labels almost
everything abnormal.

The pooled figure is therefore an average dominated by one site.
\texttt{training-e} contributes 321 of the 484 test recordings
(\SI{66}{\percent}) and has a low abnormal prevalence of \SI{8.7}{\percent}; it
is both the largest and the easiest partition. \texttt{training-a}, by contrast,
has \SI{72.1}{\percent} abnormal prevalence, and a model that answers
``abnormal'' there is right \SI{72}{\percent} of the time by accident of
prevalence.

This is exactly the behaviour Section~\ref{sec:clustering} predicted from the
unsupervised geometry, and it is what a site-prevalence shortcut looks like when
measured. We are careful about the inference: the per-site breakdown establishes
that performance does not transfer across acquisition conditions, which is a
measurement. That the mechanism is prevalence-driven site recognition is a
further inference, supported by the clustering result and by the direction of
the failures. The natural confirming experiment --- training a classifier to
predict the sub-database from the same features, and re-evaluating the diagnosis
across sites rather than pooled --- is reported next in
Section~\ref{sec:siteshortcut}.

The smaller sub-databases (\texttt{training-d}, $n=8$; \texttt{training-f},
$n=17$) carry very little weight, and their figures should be read as tendencies
rather than findings. \texttt{training-c} ($n=4$) contains a single class in the
test split and is omitted.

\subsection{Isolating the site shortcut}
\label{sec:siteshortcut}

\begin{table}[t]
\centering
\caption{Isolating the site shortcut. All three blocks are computed on the
\emph{training} partition by grouped cross-validation; the test split of
Table~\ref{tab:results} is never opened here. A single fixed, balanced random
forest is used throughout, because the quantity of interest is the gap between
evaluation protocols rather than any absolute score. Site chance is \num{0.167}
(six sites); MAcc chance is \num{0.500}.}
\label{tab:siteshortcut}
\begin{tabular}{lrr}
\toprule
 & MFCC & PWP \\
\midrule
\multicolumn{3}{l}{\emph{A. Is the site recoverable from the features?}}\\
\quad Site MAcc, segment level      & 0.647 & 0.549 \\
\quad Site accuracy, recording level & \textbf{0.871} & \textbf{0.817} \\
\midrule
\multicolumn{3}{l}{\emph{B. Does the diagnosis transfer across sites?}}\\
\quad Diagnosis MAcc, pooled CV          & 0.840 & 0.810 \\
\quad Diagnosis MAcc, leave-one-site-out & 0.457 & 0.472 \\
\quad Drop (pooled $-$ LOSO)             & \textbf{0.382} & \textbf{0.338} \\
\midrule
\multicolumn{3}{l}{\emph{C. Does signal survive once site is fixed?}}\\
\quad Diagnosis MAcc, within-site & 0.756 & 0.703 \\
\bottomrule
\end{tabular}
\end{table}


\begin{figure}[t]
\centering
\includegraphics[width=\columnwidth]{fig_site_shortcut}
\caption{Isolating the site shortcut on the training partition. Recording site
is recoverable from the features well above chance (blue); pooled
cross-validated diagnosis MAcc (orange) collapses to chance under
leave-one-site-out evaluation (green); a weak diagnosis signal nonetheless
survives once site is held fixed (red). Test data is not used.}
\label{fig:siteshortcut}
\end{figure}

The per-site breakdown shows that performance does not transfer, but it does not
by itself prove that the site is what the features encode. We test that directly,
on the training split only so that the test partition of
Table~\ref{tab:results} stays sealed, with three grouped cross-validation
experiments driven by one fixed balanced random forest
(Table~\ref{tab:siteshortcut}, Fig.~\ref{fig:siteshortcut}).

First, the site is plainly recoverable. A classifier asked to name the
sub-database from the MFCC features alone reaches \num{0.871} accuracy at the
recording level against a \num{0.167} chance rate, and \num{0.817} on PWP. The
substrate for a shortcut is therefore present in the representation itself,
before any diagnostic label is introduced.

Second, the diagnosis does not transfer across that structure. When the same
features are used for the intended task but the model is evaluated
\emph{leave-one-site-out} --- trained on five sites and tested on the held-out
sixth --- MAcc falls from \num{0.840} pooled to \num{0.457} on MFCC, and from
\num{0.810} to \num{0.472} on PWP. Both land at the \num{0.500} of a constant
predictor. The pooled figure that anchors the leaderboard describes performance
\emph{within} the corpus's site mixture; almost none of it survives a change of
acquisition site, which is the condition a screening deployment would actually
face.

Third --- and this is the qualification that keeps the conclusion honest --- the
features are not \emph{only} a site fingerprint. Held at a fixed site, where
prevalence can no longer stand in for the label, diagnosis MAcc is \num{0.756}
on MFCC and \num{0.703} on PWP, clearly above chance. That average is, however,
inflated by two untrustworthy partitions: \texttt{training-c} ($n=22$) reports a
within-site MAcc of \num{1.000} and \texttt{training-e} (abnormal prevalence
\SI{7.4}{\percent}) of \num{0.985}, both artefacts of tiny or extremely skewed
samples. On the three large, better-balanced sites the within-site MAcc is a
more sober \num{0.55}--\num{0.64}. The reading we take from all three blocks is
that the pooled metric is substantially, but not entirely, a site-prevalence
artefact: a genuine but weak diagnostic signal exists, and the aggregate score
overstates it by roughly the \num{0.34}--\num{0.38} that evaporates under
leave-one-site-out evaluation.

\subsection{Overfitting and calibration}

Two weaknesses are worth stating. Both random forests show a large gap between
training and cross-validated balanced accuracy --- \num{0.947} versus
\num{0.867} on MFCC, and \num{0.945} versus \num{0.814} on PWP --- so they are
memorising the training folds despite the depth and leaf-size grid; stronger
regularisation would be the obvious next step. Second, the classical models are
poorly calibrated, with expected calibration error \num{0.129} for RF--MFCC and
\num{0.150} for RF--PWP. Since recording-level decisions are formed by averaging
window probabilities, miscalibration propagates into the aggregation, and the
close agreement between mean-probability and majority-vote aggregation
(within \num{0.004} MAcc) suggests the probabilities are contributing little
beyond their ordering.

\subsection{Comparison with published work}

\begin{table}[t]
\centering
\caption{Comparison with published PhysioNet/CinC 2016 entries.}
\label{tab:literature}
\begin{tabular}{lrrrc}
\toprule
System & MAcc & Se. & Sp. & Comparable \\
\midrule
This work (best) & 0.871 & 0.889 & 0.853 & \checkmark \\
Potes et al. 2016 (CinC winner) & 0.860 & 0.942 & 0.778 & -- \\
Rubin et al. 2016 (CNN) & 0.839 & 0.727 & 0.951 & -- \\
Zabihi et al. 2016 (ensemble) & 0.859 & 0.869 & 0.849 & -- \\
\bottomrule
\end{tabular}
\begin{flushleft}\footnotesize Published entries were scored on the official hidden test set, which was never released. Our numbers come from a held-out split of the public training data, so the comparison is contextual rather than like-for-like.\end{flushleft}
\end{table}


Table~\ref{tab:literature} places our best result alongside published CinC 2016
entries. Our MAcc of \num{0.866} sits at the top of the \num{0.84}--\num{0.86}
band those entries occupy. This must not be read as an improvement. The published
figures were obtained on the official hidden test set, which was never released;
ours come from a self-defined split of the public training data, with a
different class balance, a different site mix and a different difficulty. The
numbers are of the same order, which is the only conclusion the comparison
supports, and it is reassuring mainly as evidence that the pipeline is not
broken.
